\chapter{Metodología experimental}
\label{cap:metodologia}

\section{Datos y selección de registros}

Se usó la base CHB-MIT Scalp EEG Database, versión 1.0.0 \parencite{shoeb2009,goldberger2000}. Se descargaron todos los registros con crisis anotadas y tres registros sin crisis por caso, elegidos al azar con semilla fija, verificando cada archivo con su suma SHA-256. Las etiquetas se tomaron de los archivos \texttt{chbXX-summary.txt}. Se detectó una inconsistencia en los metadatos de la base: \texttt{RECORDS-WITH-SEIZURES} lista \texttt{chb07\_18.edf}, pero el resumen ubica la crisis en \texttt{chb07\_19.edf}; se usó la unión de ambas fuentes para la descarga y el resumen para las etiquetas.

\pendiente{tabla de sujetos (sección 1 del notebook): horas, crisis, duración.}

\section{Preprocesamiento}

\begin{enumerate}
  \item Selección de los 18 canales bipolares del montaje ``doble banana'' presentes en todos los registros; se descartan los archivos que no los tienen.
  \item Filtro Butterworth pasabanda de orden 4 entre 0,5 y 45~Hz y filtro \emph{notch} en 60~Hz ($Q=30$), ambos aplicados hacia adelante y hacia atrás (fase cero).
  \item Ventanas de 4~s sin solapamiento. Una ventana se etiqueta como crisis si al menos el 50~\% de su duración cae dentro de una crisis anotada.
\end{enumerate}

\section{Características}

Por canal: log-potencia y potencia relativa en 5 bandas, entropía espectral normalizada, frecuencia dominante, frecuencia de borde al 90~\%, \emph{line length} y log-varianza (15 características). Por ventana: coherencia media entre los 153 pares de canales en cada banda (5). Total: $18\times 15 + 5 = 275$ características, estandarizadas con media y desvío del conjunto de entrenamiento.

\section{Partición}

\pendiente{sujetos en entrenamiento / validación / prueba y cantidad de ventanas (sección 5 del notebook).}

\section{Modelos y entrenamiento}

\begin{itemize}
  \item Regresión logística y perceptrón multicapa de dos capas ocultas (128 y 64 unidades) con \emph{dropout} 0,2.
  \item Entropía cruzada binaria ponderada por clase, minilotes de 256, hasta 40 épocas con \emph{early stopping} (paciencia 8) sobre el AUC de validación.
  \item Para L-BFGS cada época es un paso sobre el conjunto completo con hasta 10 iteraciones y búsqueda lineal de Wolfe fuerte.
\end{itemize}

\section{Diseño de los experimentos}

\begin{table}[ht]
\centering
\caption{Experimentos realizados.}
\label{tab:experimentos}
\begin{tabularx}{\linewidth}{@{}lXl@{}}
\toprule
Experimento & Factores & Hipótesis \\
\midrule
Convexo & 6 optimizadores + Nelder-Mead sobre regresión logística $L_2$; brecha a $f^\star$ & \ref{h:convexo} \\
Ajuste de tasa & 3--4 tasas por optimizador, red ReLU, semilla 0 & --- \\
Grilla & 6 optimizadores $\times$ 5 activaciones $\times$ 5 semillas = 150 & \ref{h:noconvexo}, \ref{h:adaptativos} \\
Logística & 6 optimizadores $\times$ 5 semillas, sin capas ocultas & \ref{h:noconvexo} \\
Profundidad & 2, 4, 8 capas $\times$ 4 activaciones $\times$ 3 semillas, SGD+Nesterov & \ref{h:profundidad} \\
Entropía & ventanas normales vs. crisis, por ventana y por sujeto & \ref{h:entropia} \\
\bottomrule
\end{tabularx}
\end{table}

\section{Reproducibilidad}

Todo el código está en el notebook \texttt{TP\_Final\_Crisis\_EEG.ipynb} y en los módulos de \texttt{src/}. Las semillas están fijadas, cada entrenamiento se guarda en un archivo de resultados y los experimentos se pueden interrumpir y reanudar. El entorno usado fue Python 3.12 con PyTorch, scikit-learn, SciPy y statsmodels.
